\Qhead{3.10}{Взаимодействие случайных волн}{Взаимодействие случайных волн. Генерация второй оптической гармоники, самофокусировка и самомодуляция частично когерентных волн. Преобразование спектров шумовых волн в нелинейных средах без дисперсии.}
\Subfile{3.10.1}{sections/03_statistics/q10/s01}

\recap{Нелинейное преобразование определяется высшими моментами поля. Гармоники связаны со свёртками спектров; уравнение для среднего поля требует статистического замыкания.}
\board{Четвёртый момент; гауссовский фактор два; свёртка спектров.}
\oral{Как флуктуации основной волны меняют среднюю вторую гармонику?}
\bait{Можно оставить уже разобранное различие быстрой и медленной реакции среды на спеклы.}
\sources{Цысарь~\cite{tsysar2}, гармоники; Биленко и соавт.~\cite{bilenko}, шум. Выводы — самостоятельное гауссовское спаривание.}
